% page 242 of the facsimile (image p-241 of the scan)
% block 162.6 x 246.3 mm ; left margin 8.0 mm
\pgc{-12.700mm}{0.000mm}{
\l{\cel{6}234}
\l{}
\l{}
\l{\sou{infern}o.I. (\sou{mitol.}) Loko subsula, habitata da la \textquotedbl{}ombri\textquotedbl{}}
\l{\cel{3}di la mortinti, e di qua un regiono (Tartaro) esas re-}
\l{\cel{3}zervita por la maligni, ed altra (Elizeo) esas la lojeyo}
\l{\cel{3}di la yusti. - II. (\sou{kristanismo}). Loko ube supliciesas la}
\l{\cel{3}damniti (opoze a purgatorio ed a paradizo). - EFIS.}
\l{}
\l{\sou{infiltrar}.\cel{2}(\sou{netrans., aden}).\cel{2}I. Insinuar su pokope aden la}
\l{\cel{3}pori, la interstici,la fenduri, di korpo solida. - II.}
\l{\cel{3}(\sou{metaf.})\cel{2}Penetrar pokope aden la spirito, aden la mento.}
\l{\cel{3}- DEFIS.}
\l{}
\l{\sou{infinita}. (\sou{matem.})\cel{2}To quon spirito ne konceptas kom limitoza;}
\l{\cel{3}to quo esas plu grandakam irga quanto evaluebla.- (\sou{anke}}
\l{\cel{3}\sou{metaf.}) - DEFIS.}
\l{}
\l{\sou{infinitezima}. (\sou{matem.})\cel{2}Qua relatas la quanti infinite mikra.}
\l{\cel{3}- EFIS.}
\l{}
\l{\sou{infinitivo}.\cel{2}(\sou{gram.})\cel{2}Verbo-modo qua expresas la ago verbala}
\l{\cel{3}en maniero nepreciza. - DEFIS.}
\l{}
\l{\sou{infirma}.\cel{2}Qua ne povas uzar (libere, normale) irga membro od}
\l{\cel{3}organo. - EFIS.}
\l{}
\l{\sou{infixo}. (\sou{gram.)}\cel{2}Elemento qua insertesas interne di la foni}
\l{\cel{3}ye qui konsistas radiko, por modifikar la senco di olca. -}
\l{\cel{3}DEFIS.}
\l{}
\l{\sou{inflar}. (\sou{trans., per ulo, ye ulo}). Distensar omna-sinse (korpo}
\l{\cel{3}elastika) per preso interna. - (\sou{anke metaf.}) EFIS.}
\l{}
\l{\sou{inflacar.}\cel{2}(\sou{netrans.}) (\sou{financo}). Emisar exajeroze moneto-papero.}
\l{\cel{3}- DEFIS.}
\l{}
\l{\sou{inflamar}. (\sou{trans}.)Produktar ta estado di maladeso,quan karak-}
\l{\cel{3}terizas varmeso, doloro ed (ulakaze) la tumefakteso di la}
\l{\cel{3}parto malada. - DEFIS.}
\l{}
\l{\sou{inflexar}.\cel{2}(\sou{trans}.) I. (\sou{gram.}) Modifikar la pronunco-valoro}
\l{\cel{3}di vokalo, influe da altra vokalo qua esas en la silabo}
\l{\cel{3}qua sequas nemediate. - II. (\sou{geom.})\cel{2}(\sou{arko inflexita}: qua}
\l{\cel{3}konsistas ye du taloni, tangenta per olia somiti (quale la}
\l{\cel{3}embracilo tipografiala). - EFIS.}
\l{}
\l{\sou{infloresenco.}\cel{2}(\sou{bot.})\cel{2}Dispozeso di la floro sur la stipito.}
\l{\cel{3}- EFIS.}
\l{}
\l{\sou{influar}. (\sou{trans., per, pro, pri})\cel{2}(aludante kauzo nemateriala).}
\l{\cel{3}Efikar ad ulu, ad ulo, tale ke ica modifikesas. - DEFIS.}
\l{}
\l{\sou{influenzo}.\cel{2}(patol.) Gripo tre grava.- DEFIRS.}
\l{}
\l{\sou{influxo}. (\sou{biol.}) Forco, di qua la naturo esas nekonocata, qua}
\l{\cel{3}cirkulas alonge la nervi. - DEFIS.}
}
